% !TEX program = xelatex
% Компіляція: xelatex kursova.tex && xelatex kursova.tex
\documentclass[14pt,a4paper,openany]{extreport}

\usepackage{fontspec}
\setmainfont{Times New Roman}
\newfontfamily\cyrillicfonttt{Courier New}[Script=Cyrillic]
\setmonofont{Courier New}[Script=Cyrillic]

\usepackage{polyglossia}
\setdefaultlanguage{ukrainian}

\usepackage{geometry}
\geometry{
  a4paper,
  left=25mm,
  right=10mm,
  top=20mm,
  bottom=20mm,
}

\usepackage{setspace}
\onehalfspacing
\setlength{\parindent}{1.27cm}

\usepackage{amsmath,amssymb}
\usepackage{graphicx}
\usepackage{booktabs}
\usepackage{array}
\usepackage{enumitem}
\usepackage{listings}
\usepackage{caption}
\usepackage{indentfirst}
\usepackage{etoolbox}
\usepackage[indentafter]{titlesec}
\usepackage{fancyhdr}
\usepackage{hyperref}
\usepackage{xurl}
\usepackage{bookmark}
\usepackage{lastpage}
\usepackage{totcount}
\setlength{\headheight}{17pt}

\regtotcounter{figure}
\regtotcounter{table}

\hypersetup{colorlinks=true, linkcolor=black, urlcolor=blue, citecolor=black}
\Urlmuskip=0mu plus 2mu
\emergencystretch=2em

% --- Нумерація сторінок (верхній правий кут) ---
\pagestyle{fancy}
\fancyhf{}
\fancyhead[R]{\thepage}
\renewcommand{\headrulewidth}{0pt}
\fancypagestyle{plain}{\fancyhf{}\fancyhead[R]{\thepage}}

% --- Заголовки розділів (п. 3.2, 3.4 методички) ---
\renewcommand{\chaptername}{}
\renewcommand{\thechapter}{\arabic{chapter}}

% Формат: «1 НАЗВА РОЗДІЛУ» в одному рядку (п. 3.4 методички)
\titleformat{\chapter}[block]
  {\normalfont\bfseries\centering}
  {\thechapter\ }
  {0pt}
  {\MakeUppercase}
\titlespacing*{\chapter}{0pt}{0pt}{2\baselineskip}

\titleformat{\section}
  {\normalfont\bfseries}
  {\thesection\ } % пробіл після номера, без крапки
  {0pt}
  {}

\titleformat{\subsection}
  {\normalfont\bfseries}
  {\thesubsection\ }
  {0pt}
  {}

\titlespacing*{\section}{\parindent}{2\baselineskip}{2\baselineskip}
\titlespacing*{\subsection}{\parindent}{2\baselineskip}{2\baselineskip}

% --- Підписи до рисунків і таблиць (п. 3.7, 3.8) ---
\DeclareCaptionLabelFormat{ukrainian}{#1~#2}
\captionsetup{
  labelformat=ukrainian,
  labelsep=endash,
  font=normalsize,
  skip=0.5em,
}
\captionsetup[table]{position=top, justification=justified}
\captionsetup[figure]{justification=justified, singlelinecheck=false}

% --- Структурна частина (РЕФЕРАТ, ВСТУП, …) ---
\newcommand{\structpart}[1]{%
  \clearpage
  \chapter*{\bfseries\MakeUppercase{#1}}
  \addcontentsline{toc}{chapter}{#1}%
}

% --- Заголовок абзацу у вступі (додаток Г); текст після нього — новий абзац з відступом ---
\newcommand{\introheading}[1]{%
  \par\addvspace{0.5\baselineskip}%
  \noindent\textbf{#1}\par
}

\lstset{
  basicstyle=\small\ttfamily,
  breaklines=true,
  frame=single,
  numbers=left,
  numberstyle=\tiny,
  tabsize=2,
}

% --- Рисунки: thesis/images/fig-XX-....png (див. FIGURES.md) ---
\graphicspath{{images/}}

% Показує зображення, якщо файл є; інакше — рамку-плейсхолдер
\newcommand{\maybeincludegraphics}[2][0.92\textwidth]{%
  \IfFileExists{images/#2}{%
    \includegraphics[width=#1]{#2}%
  }{%
    \fbox{\parbox[c][6cm][c]{#1}{%
      \centering\small
      \textbf{Місце для рисунка}\\[0.8em]
      \texttt{thesis/images/#2}\\[0.5em]
      Опис: \texttt{FIGURES.md}%
    }}%
  }%
}

% ========== ДАНІ ВИКОНАВЦЯ (замініть на свої) ==========
\newcommand{\WorkTitle}{%
  МОНОКУЛЯРНА ВІЗУАЛЬНА ОДОМЕТРІЯ \\
  В РЕАЛЬНОМ ЧАСІ ДЛЯ БПЛА}
\newcommand{\StudentCourse}{3-го}
\newcommand{\StudentGender}{виконала}
\newcommand{\StudentGenderShort}{студентка}
\newcommand{\StudentGenderShortUpper}{Студентка}
\newcommand{\StudentName}{Уточкіна Яна Максимівна}
\newcommand{\StudentProgram}{ОПП «Інформатика»}
\newcommand{\SupervisorTitle}{%
  завідувач кафедри теоретичної кібернетики\\
  доктор фізико-математичних наук}
\newcommand{\SupervisorName}{Терещенко Василь Миколайович}
\newcommand{\Year}{2026}

\begin{document}

% =============================================================================
% Титульний аркуш (додаток А — зразок курсової роботи)
% =============================================================================
\begin{titlepage}
  \thispagestyle{empty}
  \begin{center}
    \textbf{КИЇВСЬКИЙ НАЦІОНАЛЬНИЙ УНІВЕРСИТЕТ}\\
    \textbf{ІМЕНІ ТАРАСА ШЕВЧЕНКА}\\[0.5em]
    Факультет комп'ютерних наук та кібернетики\\
    Кафедра математичної інформатики\\[2.5em]

    Курсова робота\\
    за спеціальністю 122 Комп'ютерні науки\\
    на тему:\\
    {\bfseries\WorkTitle}\\[2.5em]
  \end{center}

  \noindent
  \begin{tabular}{@{}ll}
    \StudentGender\ \StudentGenderShort\ \StudentCourse\ курсу & \\
    \StudentProgram \\[-0.5em]
    \StudentName & \hspace{5cm} \begin{tabular}{c} \underline{\hspace{3cm}} \\ \small(підпис) \end{tabular}
  \end{tabular}
  \\[2.5em]

  \noindent
  Науковий керівник:\\
  \begin{tabular}{@{}ll}
    \SupervisorTitle \\[-0.5em]
    \SupervisorName & \hspace{2cm} \begin{tabular}{c} \underline{\hspace{3cm}} \\ \small(підпис) \end{tabular}
  \end{tabular}
  \\[2.5em]

  \begin{flushright}
    Засвідчую, що в цій роботі\\
    немає запозичень з праць інших авторів\\
    без відповідних посилань.\\
    \StudentGenderShortUpper\hspace{2em}\underline{\hspace{3cm}}\\[0.2em]
    (підпис)
  \end{flushright}
  \vspace{3em}

  \begin{center}
    КИЇВ~\Year
  \end{center}
\end{titlepage}

% =============================================================================
% РЕФЕРАТ (додаток Б)
% =============================================================================
\structpart{Реферат}

Обсяг роботи \pageref{LastPage}~сторінок,
\total{figure}~ілюстрацій,
\total{table}~таблиці,
9~джерел посилань.

\smallskip
\noindent\textbf{ВІЗУАЛЬНА ОДОМЕТРІЯ, КАЛІБРУВАННЯ КАМЕРИ, МОНОКУЛЯРНА КАМЕРА, ОПТИЧНИЙ ПОТІК,}\\
\textbf{ВЕБ-ЗАСТОСУНОК, КОМП'ЮТЕРНИЙ ЗІР.}

\smallskip
Об'єктом роботи є процес оцінювання траєкторії руху монокулярної камери за послідовністю зображень у режимі реального часу або з відеофайлу. Предметом роботи є методи детекції та відстеження ознак, відновлення відносної пози камери та програмна реалізація веб-системи візуальної одометрії.

Метою роботи є розробка та дослідження програмного комплексу монокулярної візуальної одометрії з веб-інтерфейсом для візуалізації тривимірної траєкторії.

Методи розроблення: аналіз літератури з візуальної одометрії, проєктування клієнт--серверної архітектури, реалізація алгоритмів комп'ютерного зору (OpenCV), компонентна веб-розробка (Angular), експериментальна перевірка на наборах даних EuRoC MAV. Інструменти: Python, FastAPI, OpenCV, Angular, Plotly.js, Docker.

Результати роботи: спроєктовано та реалізовано веб-систему з транспортом WebSocket; розроблено модуль VO на базі Shi--Tomasi, Lucas--Kanade та суттєвої матриці з RANSAC; реалізовано калібрування камери за шаховою дошкою; отримано профіль камери з похибкою репроєкції близько 1{,}21~px на даних EuRoC; реалізовано експорт траєкторії у формат TUM та кількісну оцінку глобального масштабу за ground truth EuRoC (Sim(3)-вирівнювання, ATE).

Система може застосовуватися для навчальної демонстрації принципів візуальної одометрії, прототипування алгоритмів локалізації БПЛА та попередньої оцінки траєкторії за відео з бортової камери.

% =============================================================================
% ЗМІСТ
% =============================================================================


\tableofcontents

% =============================================================================
% СКОРОЧЕННЯ ТА УМОВНІ ПОЗНАКИ (додаток В)
% =============================================================================
\structpart{Скорочення та умовні познаки}

\noindent
API --- програмний інтерфейс застосунку (Application Programming Interface);\\
ATE --- абсолютна похибка траєкторії (Absolute Trajectory Error);\\
БПЛА --- безпілотний літальний апарат;\\
GNSS --- глобальна навігаційна супутникова система;\\
IMU --- інерціальний вимірювальний блок;\\
LK --- метод Lucas--Kanade;\\
RANSAC --- RANdom SAmple Consensus, оцінювання з випадковою підвибіркою;\\
RPE --- відносна похибка пози (Relative Pose Error);\\
SLAM --- одночасна локалізація та побудова карти (Simultaneous Localization and Mapping);\\
VO --- візуальна одометрія (Visual Odometry);\\
YAML --- мова серіалізації даних YAML Ain't Markup Language;\\
EMA --- експоненційне ковзне середнє (Exponential Moving Average);\\
REST --- стиль обміну даними HTTP (Representational State Transfer);\\
Sim(3) --- перетворення подібності (масштаб, поворот, зсув) у тривимірному просторі;\\
TUM --- формат файлу траєкторії (Technical University of Munich);\\
$E$ --- суттєва матриця;\\
$K$ --- матриця внутрішніх параметрів камери.

% =============================================================================
% ВСТУП (додаток Г)
% =============================================================================
\structpart{Вступ}

\introheading{Оцінка сучасного стану об'єкта розроблення.}

Визначення положення та орієнтації мобільного робота або безпілотного літального апарата є фундаментальною задачею автономної навігації. У середовищах, де сигнал GNSS недоступний або ненадійний, основним джерелом інформації про рух часто стає бортова відеокамера. Візуальна одометрія дозволяє оцінювати відносний рух камери між послідовними кадрами шляхом інтегрування зміщень та обертань у траєкторію [5]. Монокулярна VO використовує один відеопотік, що спрощує апаратну частину, але створює труднощі через неоднозначність масштабу та накопичення похибки. Сучасні SLAM-системи, зокрема ORB-SLAM [2] та PTAM [3], демонструють високу точність, однак їх повна реалізація є складною для навчальних та прототипних задач.

\introheading{Актуальність роботи та підстави для її виконання.}

Розробка веб-орієнтованої системи, яка поєднує алгоритми комп'ютерного зору на сервері та інтерактивну візуалізацію траєкторії в браузері, актуальна для навчального процесу, швидкого прототипування та попередньої оцінки руху БПЛА за відео. Наявність відкритих наборів даних, зокрема EuRoC MAV [1], дозволяє верифікувати калібрування камери та якість VO.

\introheading{Мета й завдання роботи.}

Метою курсової роботи є розробка та дослідження програмного комплексу монокулярної візуальної одометрії з веб-інтерфейсом, що працює в режимі реального часу або з попередньо записаного відео. Для досягнення мети поставлено такі завдання:

\begin{itemize}[leftmargin=\parindent, itemsep=0pt]
  \item проаналізувати теоретичні основи монокулярної VO та калібрування камери;
  \item спроєктувати клієнт--серверну архітектуру системи;
  \item реалізувати модуль оцінювання руху на основі оптичного потоку та епіполярної геометрії;
  \item реалізувати калібрування камери за зображеннями шахової дошки;
  \item забезпечити візуалізацію тривимірної траєкторії та керування профілями камери;
  \item провести експериментальну перевірку на наборах даних EuRoC MAV, зокрема кількісну оцінку траєкторії за ground truth.
\end{itemize}

\introheading{Об'єкт, методи й засоби розроблення.}

Об'єктом дослідження є процес оцінювання траєкторії руху монокулярної камери за послідовністю зображень. Предметом --- методи детекції ознак, відстеження, відновлення пози та програмна реалізація веб-системи VO. У роботі застосовано методи аналітичної геометрії, комп'ютерного зору (OpenCV), проєктування клієнт--серверних систем, компонентний підхід у веб-розробці (Angular), а також експериментальну верифікацію на еталонних наборах даних.

\introheading{Можливі сфери застосування.}

Створена система може використовуватися для демонстрації принципів VO в навчальному процесі, швидкого прототипування алгоритмів локалізації, попередньої оцінки траєкторії БПЛА та збереження параметрів калібрування камери у профілях для подальших сеансів VO.

\introheading{Структура роботи.}

Робота складається з вступу, чотирьох розділів, висновків, переліку джерел посилання та додатків. У розділі~1 викладено теоретичні основи; у розділі~2 --- проєктування системи; у розділі~3 --- реалізацію; у розділі~4 --- експериментальні дослідження.

% =============================================================================
% РОЗДІЛ 1
% =============================================================================
\chapter{Теоретичні основи монокулярної візуальної одометрії}

\section{Модель камери з центральною перспективою}

Камеру часто апроксимують центральною перспективною моделлю («камера-обскура», pin-hole). Точка світового простору $\mathbf{P} = (X, Y, Z)^\mathsf{T}$ проєктується на зображення через матрицю внутрішніх параметрів $K$:
\begin{equation}
  \lambda
  \begin{bmatrix} u \\ v \\ 1 \end{bmatrix}
  =
  K
  \begin{bmatrix} X \\ Y \\ Z \end{bmatrix},
  \qquad
  K =
  \begin{bmatrix}
    f_u & 0   & c_u \\
    0   & f_v & c_v \\
    0   & 0   & 1
  \end{bmatrix},
  \label{eq:K}
\end{equation}
\noindent
де $f_u$, $f_v$ --- фокусні відстані в пікселях; $c_u$, $c_v$ --- координати головного пункту; $\lambda$ --- масштабний множник глибини.

На рис.~\ref{fig:pinhole} схематично показано проєкцію точки світового простору на площину зображення та систему координат OpenCV.

\begin{figure}[htbp]
  \centering
  \maybeincludegraphics{fig-01-pinhole.png}
  \caption{Модель камери з центральною перспективою та система координат OpenCV}
  \label{fig:pinhole}
\end{figure}

Для точної геометрії враховують дисторсію об'єктива, яку компенсують функцією усунення дисторсії \texttt{undistort} у OpenCV [4].

\section{Калібрування камери}

Калібрування --- визначення $K$ та коефіцієнтів дисторсії за зображеннями відомого просторового шаблону. Найпоширеніший підхід --- шахова дошка (метод Чжана, функція OpenCV \texttt{calibrateCamera}):
\begin{enumerate}[leftmargin=*]
  \item для кожного кадра знаходять внутрішні кути сітки;
  \item уточнюють їх субпіксельно;
  \item мінімізують похибку репроєкції --- RMS-відстань між спостереженими та проєктованими кутами.
\end{enumerate}

Якість калібрування оцінюють за похибкою репроєкції: значення менше 0{,}5--1{,}0~px вважають прийнятними для VO.

\section{Детекція та відстеження ознак}

У розробленій системі застосовано:
\begin{itemize}[leftmargin=*]
  \item метод Shi--Tomasi (\texttt{goodFeaturesToTrack}) для вибору сильних кутів;
  \item метод Lucas--Kanade --- пірамідальний оптичний потік\\ (\texttt{calcOpticalFlowPyrLK});
  \item перевірку узгодженості прямого та зворотного потоку --- відхилення треку, якщо зворотний потік не повертає точку в початкове положення.
\end{itemize}

Приклад відстеження ознак між двома послідовними кадрами наведено на рис.~\ref{fig:flow}.

\begin{figure}[htbp]
  \centering
  \maybeincludegraphics{fig-02-optical-flow.png}
  \caption{Відстеження ознак методом Lucas--Kanade між двома кадрами}
  \label{fig:flow}
\end{figure}

\section{Епіполярна геометрія та суттєва матриця}

Для двох кадрів з відомою калібровкою $K$ відносний рух описується суттєвою матрицею $E$ [6]:
\begin{equation}
  E = [\mathbf{t}]_\times R,
  \label{eq:essential}
\end{equation}
\noindent
де $R$ --- відносне обертання, $\mathbf{t}$ --- напрямок переносу (одиниця довжини), $[\mathbf{t}]_\times$ --- кососиметрична матриця.

Матрицю $E$ оцінюють методом RANSAC (\texttt{findEssentialMat}), після чого \texttt{recoverPose} відновлює $R$ та $\mathbf{t}$ з перевіркою зодисаальності (cheirality).

\section{Неоднозначність масштабу в монокулярній VO}

На відміну від стерео- або RGB-D-систем (камера з глибиною), одна камера не спостерігає абсолютну глибину без додаткових припущень. Вектор переносу $\mathbf{t}$ відновлюється лише з точністю до множника [5]. У даній роботі застосовано спрощений підхід: нормалізація кроку переносу та обмеження відносного масштабу між кадрами. За наявності еталонної траєкторії (ground truth EuRoC) метричний масштаб відновлюється \emph{a posteriori} через Sim(3)-вирівнювання (розділ~4).

\section{Огляд підходів ORB-SLAM та PTAM}

Сучасні SLAM-системи (ORB-SLAM [2], PTAM [3]) поєднують надійне відстеження з дескрипторами, ключові кадри, локальне пакетне уточнення та замикання контуру. Розроблена система реалізує полегшений VO-конвеєр, орієнтований на обробку у реальному часі через WebSocket, без глобальної карти та замикання контурів.

% =============================================================================
% РОЗДІЛ 2
% =============================================================================
\chapter{Проєктування системи}

\section{Загальна архітектура}

Система побудована за архітектурою клієнт--сервер (рис.~\ref{fig:arch}). Обчислювально складні операції комп'ютерного зору виконуються на сервері (Python); клієнт на Angular відповідає за захоплення відео, налаштування та візуалізацію.

\begin{figure}[htbp]
  \centering
  \maybeincludegraphics{fig-03-architecture.png}
  \caption{Логічна архітектура системи}
  \label{fig:arch}
\end{figure}

\section{Режими роботи}

Підтримуються два режими: «Потік» --- захоплення з вебкамери у реальному часі; «З файлу» --- обробка завантаженого відеофайлу (наприклад, конвертованої послідовності EuRoC). Перед запуском VO користувач обирає профіль калібрування камери.

\section{Протокол обміну даними}

З'єднання встановлюється через WebSocket за адресою
\url{ws://localhost:8000/ws/vo-stream?config_id=<profile_id>}.
Цикл «запит--відповідь»: клієнт надсилає кадр (JPEG у кодуванні Base64); сервер обробляє кадр у окремому потоці (\texttt{asyncio.to\_thread}); сервер повертає координати та стан відстеження; клієнт додає точку на тривимірний графік і надсилає наступний кадр.

\section{Зберігання конфігурації}

Профілі камери зберігаються у форматі YAML у каталозі \path{vo-uav/storage/configs/}. Профіль містить внутрішні параметри $f_u$, $f_v$, $c_u$, $c_v$, коефіцієнти дисторсії, метадані калібрування та параметри VO.

\section{Стек технологій}

Стек технологій наведено в табл.~\ref{tab:stack}.

\begin{table}[htbp]
  \caption{Технологічний стек системи}
  \label{tab:stack}
  \centering
  \begin{tabular}{@{}>{\raggedright\arraybackslash}p{2.7cm}>{\raggedright\arraybackslash}p{10.3cm}@{}}
    \toprule
    Шар & Технології \\
    \midrule
    Серверна частина & Python 3.11+, FastAPI, OpenCV, NumPy, PyYAML \\
    Клієнтська частина & Angular 21, TypeScript, Angular Material, Plotly.js \\
    Транспорт & WebSocket, REST (калібрування, профілі, експорт траєкторії) \\
    Розгортання & Docker, Docker Compose \\
    \bottomrule
  \end{tabular}
\end{table}

% =============================================================================
% РОЗДІЛ 3
% =============================================================================
\chapter{Реалізація програмного комплексу}

\section{Модуль візуальної одометрії}

Клас \texttt{VisualOdometry} (\texttt{vo-uav/src/vo.py}) реалізує інкрементальний VO-конвеєр. Система координат OpenCV: $X$ --- вправо, $Y$ --- вниз, $Z$ --- углиб сцени.

\subsection{Попередня обробка}

Виконується масштабування $K$ до фактичної роздільності кадру (якщо вона відрізняється від калібрувальної) та опційне усунення дисторсії за коефіцієнтами з профілю.

\subsection{Конвеєр обробки кадру}

Для кожного кадру $I_t$:
\begin{enumerate}[leftmargin=*]
  \item відстеження ознак між $I_{t-1}$ та $I_t$ (LK + прямий--зворотний потік);
  \item оцінка відносної пози: основний шлях --- $E$ + RANSAC + \texttt{recoverPose}; резервний --- \texttt{estimateAffinePartial2D} для плоских рухів;
  \item оцінка масштабу кроку (відношення глибин triangulation або медіана оптичного потоку);
  \item інтеграція пози; поповнення ознак; періодичне оновлення ключового кадру;
  \item згладжування траєкторії (EMA + відсіювання стрибків).
\end{enumerate}

Блок-схема конвеєра наведена на рис.~\ref{fig:pipeline}.

\begin{figure}[htbp]
  \centering
  \maybeincludegraphics[0.88\textwidth]{fig-04-vo-pipeline.png}
  \caption{Блок-схема конвеєра візуальної одометрії}
  \label{fig:pipeline}
\end{figure}

\subsection{Стани відстеження}

Програмний інтерфейс повертає поле \texttt{tracking} зі значеннями \texttt{initializing}, \texttt{ok}, \texttt{degraded}, \texttt{lost} (ініціалізація, успішно, погіршено, втрачено) для діагностики якості оцінки в інтерфейсі користувача.

\section{Модуль калібрування}

Файл \texttt{camera\_calibration.py} реалізує побудову тривимірного шаблону кутів шахової дошки, детекцію з \texttt{findChessboardCornersSB} та класичним резервним методом, автоматичний підбір розміру сітки (зокрема 6$\times$7 для EuRoC \texttt{cam\_checkerboard}), виклик \texttt{calibrateCamera} та збереження результату. Кінцева точка REST: \texttt{POST /api/configs/calibrate}.

\section{Експорт траєкторії та оцінювання масштабу}

Модуль \texttt{trajectory\_eval.py} забезпечує експорт накопиченої траєкторії у форматі TUM (\texttt{timestamp tx ty tz qx qy qz qw}) та кількісне порівняння з еталоном EuRoC. Для кожного кадру VO зберігається індекс кадру; timestamp береться з \texttt{mav0/cam0/data.csv} відповідної послідовності. Ground truth інтерполюється з \texttt{state\_groundtruth\_estimate0/data.csv} на ці мітки часу.

Глобальний масштаб $s$ обчислюється алгоритмом Umeyama (Sim(3)-вирівнювання): еталонна траєкторія $\approx s \cdot R \cdot \mathbf{p}_\text{VO} + \mathbf{t}$. Після масштабування обчислюється абсолютна похибка траєкторії ATE. REST-кінцеві точки: \texttt{GET /api/trajectory/euroc-sequences}, \texttt{POST /api/trajectory/export-tum}, \texttt{POST /api/trajectory/evaluate}. У вкладці «З файлу» доступні кнопки \emph{Export TUM}, \emph{Export scaled TUM} та \emph{Compute global scale}. Для пакетної оцінки використовується скрипт \texttt{scripts/run\_vo\_eval.py}.

\section{Клієнтська частина та розгортання}

Angular-додаток містить вкладки «Потік» / «З файлу», елементи керування профілем камери, робочі простори VO з тривимірним графіком Plotly, вибір послідовності EuRoC для порівняння з еталоном та сервіс WebSocket зі сторожовим механізмом (watchdog). Головний інтерфейс показано на рис.~\ref{fig:ui}. Каталог \texttt{Datasets} змонтовано в контейнер backend (\texttt{EUROC\_DATASETS\_PATH=/datasets}). Розгортання здійснюється через \texttt{docker-compose.yml} (сервер --- порт 8000, клієнт --- порт 4200).

\begin{figure}[htbp]
  \centering
  \maybeincludegraphics{fig-05-ui.png}
  \caption{Головний інтерфейс системи (вкладка «Потік»)}
  \label{fig:ui}
\end{figure}

% =============================================================================
% РОЗДІЛ 4
% =============================================================================
\chapter{Експериментальні дослідження}

\section{Набір даних EuRoC MAV}

Для тестування використано набір EuRoC MAV Dataset (ETH Zurich) [1] --- стандартний еталон для візуально-інерційної одометрії та SLAM. Послідовність \texttt{cam\_checkerboard} містить кадри камери \texttt{cam0} (752$\times$480~px, 20~Гц) з видимою шаховою дошкою. Для калібрування потрібна сітка 6$\times$7 внутрішніх кутів. З 2376~кадрів дошку успішно детектовано на 1740.

\section{Результати калібрування}

Параметри профілю dataset.yaml, отриманого за 11~зображеннями з EuRoC, наведено в табл.~\ref{tab:calib}.

\begin{table}[htbp]
  \caption{Параметри калібрування (профіль Dataset)}
  \label{tab:calib}
  \centering
  \begin{tabular}{@{}lr@{}}
    \toprule
    Параметр & Значення \\
    \midrule
    $f_u$ & 397,88~px \\
    $f_v$ & 394,66~px \\
    $c_u$ & 360,80~px \\
    $c_v$ & 271,66~px \\
    Похибка репроєкції & 1,21~px \\
    Роздільність & 752 $\times$ 480~px \\
    Внутрішні кути & 6 $\times$ 7 \\
    Розмір клітини & 38~mm \\
    \bottomrule
  \end{tabular}
\end{table}

Похибка репроєкції близько 1,2~px --- прийнятний рівень; для підвищення точності VO рекомендується більше різноманітних ракурсів. Приклад детекції кутів шахової дошки наведено на рис.~\ref{fig:calib}.

\begin{figure}[htbp]
  \centering
  \maybeincludegraphics{fig-06-calibration.png}
  \caption{Детекція внутрішніх кутів шахової дошки 6×7\\ (EuRoC \texttt{cam\_checkerboard})}
  \label{fig:calib}
\end{figure}

\section{Оцінювання VO}

Якісна оцінка включає аналіз форми траєкторії на Plotly (рис.~\ref{fig:traj}), полів \texttt{confidence} (впевненість) та \texttt{tracking} (стан відстеження), поведінку при нерухомій камері та порівняння типового профілю EuRoC і власної калібровки.

\begin{figure}[htbp]
  \centering
  \maybeincludegraphics{fig-07-trajectory.png}
  \caption{Тривимірна траєкторія, побудована системою VO\\ (графік Plotly scatter3d)}
  \label{fig:traj}
\end{figure}

Для кількісної оцінки на польових послідовностях EuRoC Vicon Room використано еталонну траєкторію (ground truth) з файлу \texttt{state\_groundtruth\_estimate0/data.csv}. Алгоритм оцінювання:
\begin{enumerate}[leftmargin=\parindent, itemsep=0pt]
  \item експорт траєкторії VO у формат TUM;
  \item інтерполяція ground truth на timestamp камери \texttt{cam0};
  \item Sim(3)-вирівнювання (Umeyama) та обчислення глобального масштабу $s$;
  \item обчислення ATE (RMSE, median, max) між масштабованою та еталонною траєкторіями.
\end{enumerate}

Експеримент проведено на відео \texttt{V1\_01\_easy.mp4} (2912~кадрів, 752$\times$480~px, 20~Гц), отриманому з кадрів \texttt{cam0} EuRoC, з профілем \texttt{euroc\_default}. Результати наведено в табл.~\ref{tab:vo-eval}.

\begin{table}[htbp]
  \caption{Кількісна оцінка VO на послідовності EuRoC \texttt{V1\_01\_easy}}
  \label{tab:vo-eval}
  \centering
  \begin{tabular}{@{}lr@{}}
    \toprule
    Метрика & Значення \\
    \midrule
    Кадрів у стані \texttt{ok} & 94,9\,\% \\
    Середня \texttt{confidence} & 0,957 \\
    Глобальний масштаб $s$ (Sim(3)) & 0,063 \\
    ATE RMSE & 1,72~м \\
    ATE median & 1,52~м \\
    Довжина шляху GT & 58,6~м \\
    Довжина шляху VO (масштабована) & 9,2~м \\
    \bottomrule
  \end{tabular}
\end{table}

Високий відсоток кадрів у стані \texttt{ok} свідчить про стабільність відстеження ознак. ATE порядку 1,7~м на шляху $\approx$59~м є типовим для монокулярного VO без loop closure. Суттєве недооцінювання довжини масштабованої траєкторії (9,2~м проти 58,6~м) підтверджує накопичення помилки масштабу та дрейф форми траєкторії --- обмеження, притаманні монокулярному VO.

Додатково траєкторії можна порівнювати з evo [7] (RPE, детальні графіки помилок) за експортованими TUM-файлами.

\section{Обмеження експериментів}

Монокулярна VO без замикання контуру накопичує дрейф; масштаб траєкторії під час польоту залишається умовним --- метричний масштаб відновлюється лише \emph{a posteriori} за наявності ground truth; сценарії з вебкамерою часто активують резервний affine-метод; затримка мережі WebSocket впливає на частоту кадрів у режимі «Потік».

% =============================================================================
% ВИСНОВКИ
% =============================================================================
\structpart{Висновки}

У курсовій роботі спроєктовано та реалізовано веб-систему монокулярної візуальної одометрії з клієнт--серверною архітектурою, транспортом WebSocket і тривимірною візуалізацією траєкторії.

Розроблено VO-модуль на базі Shi--Tomasi, Lucas--Kanade, суттєвої матриці з RANSAC, резервного affine-методу та післяобробки згладжування. Реалізовано калібрування камери за шаховою дошкою з автопідбором розміру сітки та збереженням профілів YAML.

На даних EuRoC \texttt{cam\_checkerboard} отримано профіль з похибкою репроєкції близько 1,21~px; підтверджено коректність детекції дошки 6$\times$7. На послідовності \texttt{V1\_01\_easy} реалізовано експорт траєкторії у формат TUM, обчислення глобального масштабу $s \approx 0{,}063$ (Sim(3)) та ATE RMSE $\approx 1{,}72$~м; стабільність відстеження --- 94,9\,\% кадрів у стані \texttt{ok}.

Система придатна для навчальної демонстрації VO та прототипування. Перспективи розвитку: метрики RPE через evo, локальне пакетне уточнення, ORB-дескриптори, поєднання з IMU, замикання контуру.

% =============================================================================
% ПЕРЕЛІК ДЖЕРЕЛ ПОСИЛАННЯ (ДСТУ 8302:2015, додаток Д)
% =============================================================================
\structpart{Перелік джерел посилання}

\begin{enumerate}[label=\arabic*., leftmargin=*, itemsep=0.35em]

\item Burri M., Nikolic J., Gohl P. [та ін.]. The EuRoC micro aerial vehicle datasets. International Journal of Robotics Research. 2016.

\item Mur-Artal R., Montiel J., Tardós J.~D. ORB-SLAM: a versatile and accurate monocular SLAM system. IEEE Transactions on Robotics. 2015.

\item Klein G., Murray D. Parallel tracking and mapping for small AR workspaces. ISMAR. 2007.

\item Camera calibration with OpenCV // OpenCV Documentation. URL: \url{https://docs.opencv.org/4.x/d4/d94/tutorial_camera_calibration.html} (дата звернення: 29.05.2026).

\item Scaramuzza D., Fraundorfer F. Visual Odometry: Part I --- The First 30 Years and Fundamentals. IEEE Robotics \& Automation Magazine. 2011.

\item Hartley R., Zisserman A. Multiple View Geometry in Computer Vision. Cambridge University Press, 2004. 655~с.

\item evo: Python package for the evaluation of odometry and SLAM // GitHub. URL: \url{https://github.com/MichaelGrupp/evo} (дата звернення: 29.05.2026).

\item FastAPI Documentation. URL: \url{https://fastapi.tiangolo.com/} (дата звернення: 29.05.2026).

\item Angular Documentation. URL: \url{https://angular.dev/} (дата звернення: 29.05.2026).

\end{enumerate}

\end{document}
